\documentclass{article}

% Packages
\usepackage[margin=1in]{geometry}
\usepackage{amsmath,amssymb,amsthm}
\usepackage[draft]{graphicx}
\usepackage{booktabs}
\usepackage{hyperref}
\usepackage{xcolor}
\usepackage{subcaption}
\usepackage{float}
\usepackage{multirow}
\usepackage{microtype}
\usepackage{natbib}
\usepackage{url}

% Convenience macros
\newcommand{\ie}{i.e.,\xspace}
\newcommand{\eg}{e.g.,\xspace}
\newcommand{\bary}{\mathbf{b}}
\newcommand{\R}{\mathbb{R}}
\newcommand{\Kk}{K}

\title{Finding Belief Geometries with Sparse Autoencoders}

\author{
  [Author Names Omitted for Review]
}

\date{}

\begin{document}

\maketitle

\begin{abstract}
Understanding the geometric structure of internal representations is a central goal of mechanistic interpretability.
Prior work has shown that transformers trained on sequences generated by hidden Markov models can encode probabilistic
belief states as simplex-shaped geometries in their residual stream representations, with vertices corresponding to
latent generative states. Whether large language models trained on naturalistic text develop analogous geometric
representations remains an open question.

We introduce a pipeline for discovering simplex-structured subspaces in large language models that combines sparse
autoencoders (SAEs), $k$-subspace latent clustering, and a simplex-fitting autoencoder (AANet). Our contributions are:
(1) a general pipeline for identifying candidate simplex geometries in transformer representations, (2) validation of
the approach on a controlled transformer trained on a multipartite hidden Markov model with known belief-state
geometry, and (3) empirical evidence for candidate simplex structures in the SAE latent space of Gemma-2-9B. As an
initial validation, we apply the pipeline to a small transformer trained on a multipartite hidden Markov model with
known ground-truth belief geometries and show that the method recovers the expected simplex structure. We then apply
the pipeline to the SAE latent space of Gemma-2-9B, searching for simplex geometries with $K \ge 3$, and identify 13
priority clusters exhibiting strong candidate simplex structure; all selected clusters correspond to 2-simplices
($K = 3$).

We evaluate these clusters using four complementary tests: cross-vertex KL divergence of next-token distributions,
the predictive advantage of barycentric simplex coordinates over individual latent features, causal steering of model
continuations via AANet-guided activations, and semantic coherence of vertex labels produced by a frontier language
model. For 3 of the 13 clusters, barycentric simplex coordinates outperform the best individual latent at predicting
next-token distributions (Wilcoxon $p < 10^{-14}$), an advantage not observed in size-matched null clusters. Causal
interventions guided by simplex coordinates shift model continuations toward target vertex semantics in 8 of 13
clusters, with the strongest cluster achieving a mean shift rate of 0.42.

These findings provide preliminary empirical evidence that mixtures over latent generative states may be encoded as
simplex-like geometries within subspaces of the Gemma-2-9B residual representation.
\end{abstract}
  
% ============================================================
\section{Introduction}
\label{sec:intro}
% ============================================================

Understanding the geometric structure of internal representations is a central goal of mechanistic interpretability.  A growing body of work
suggests that transformer language models organize their activations into structured geometric objects, including sparse feature
directions, low-dimensional subspaces, and curved manifolds that encode latent variables relevant to the model's predictions
\citep{elhage2022toy, cunningham2023sparse, bricken2023monosemanticity, gurnee2023language}.  Interpreting these geometric structures provides
a promising route to understanding how large language models represent information and perform computation.

Many recent discoveries of geometric structure in transformer representations have arisen in relatively constrained algorithmic
settings.  For example, models have been shown to encode numerical quantities such as token position, character counts, or integer values
on low-dimensional manifolds that support algorithmic operations \citep{gurnee2025manifolds, kantamneni2025trigonometry}.  In these
cases, the learned geometry often reflects the underlying algebraic structure of the task: integer values may lie on circular or helical
manifolds that enable modular arithmetic, while counting variables may be embedded in curved manifolds manipulated by attention heads to
compute distances or offsets.

While these results demonstrate that transformers can exploit structured geometry when learning algorithmic tasks, it remains less
clear whether similar geometric representations arise when models track more abstract latent variables during natural language
processing.  In many real-world generative processes, the next token depends on an unobserved latent state that evolves over time.  One
principled hypothesis is that effective sequence predictors internally track probabilistic beliefs over such latent generative states.

Recent work has shown that transformers trained on synthetic sequences generated by hidden Markov models can represent these belief
distributions as simplex-shaped geometries in their residual stream activations, with vertices corresponding to the latent states of the
underlying process \citep{shai2024transformers, piotrowski2025constrained}. In these models, the internal representation of uncertainty over latent
states forms a geometric simplex whose barycentric coordinates encode the model's belief distribution.

Whether large language models trained on naturalistic text develop analogous geometric representations remains an open question.  Unlike
synthetic HMM settings, the latent variables relevant for natural language---such as discourse states, semantic frames, or topic
structure---are complex, weakly supervised, and not directly observable.  Identifying geometric structures that encode such latent
variables therefore requires methods that can search automatically for candidate geometries in large models.

In this paper we introduce a pipeline for discovering candidate simplex structures in transformer representations.  The method
combines sparse autoencoders (SAEs), which provide an interpretable feature basis for the residual stream, with $k$-subspace clustering to
identify groups of co-activating latent features.  These candidate subspaces are then analyzed using a simplex-fitting autoencoder
(AANet) to test whether the observed activations are well-described by low-dimensional simplex geometries.

We first validate the pipeline on a controlled toy model: a small transformer trained on a multipartite hidden Markov model with known
ground-truth belief geometry.  In this setting, the method reliably recovers the expected simplex representations.  We then apply the
pipeline to the SAE latent space of Gemma-2-9B, searching for simplex geometries with $K \ge 3$.  From this search we identify 13 priority
clusters exhibiting strong candidate simplex structure; all selected clusters correspond to 2-simplices ($K = 3$).

We evaluate these candidate structures using four independent validation strands:
\begin{enumerate}
\item \textbf{Cross-vertex predictive differences.} We measure differences in next-token
prediction distributions across simplex vertices.

\item \textbf{Barycentric predictive advantage.} We test whether barycentric simplex
coordinates predict next-token distributions better than any individual latent feature.

\item \textbf{Causal steering.} We intervene on model activations using AANet-guided
simplex coordinates and measure whether model continuations shift toward the
semantics associated with the target vertex.

\item \textbf{Semantic vertex coherence.} We evaluate whether tokens near each vertex
share coherent semantic structure, using a frontier language model to generate
vertex descriptions.
\end{enumerate}

These tests probe predictive, causal, and semantic implications of the simplex hypothesis. Taken together, the results
provide suggestive evidence that some subspaces of the SAE latent representation in Gemma-2-9B behave like simplex
structures encoding mixtures over latent generative states. More broadly, our findings indicate that the geometric
representations previously observed in controlled algorithmic settings may also arise in the context of more abstract
latent variables involved in natural language processing.

% ============================================================
\section{Related Work and Background}
\label{sec:background}
% ============================================================

\subsection{Sparse feature representations in transformers}

A central line of research in mechanistic interpretability aims to decompose transformer activations into interpretable feature
directions.  Sparse autoencoders (SAEs) applied to the residual stream have been shown to recover large collections of sparse features that
correspond to human-interpretable concepts and linguistic structures \citep{cunningham2023sparse, bricken2023monosemanticity}.  These
features provide a useful basis for studying the internal structure of transformer representations, as many higher-level behaviors can be
understood as interactions among sparse feature activations.

Our work builds on this paradigm by analyzing geometric structure in the space of SAE latent activations.  Rather than focusing on
individual feature directions, we search for low-dimensional geometric objects formed by coordinated patterns of latent feature activity.

\subsection{Geometric structure in transformer representations}

Beyond individual feature directions, recent work has identified structured geometric objects in transformer representations.  For
example, linear subspaces have been shown to encode grammatical and semantic variables \citep{gurnee2023language}.  Other studies have
found that model activations can lie on low-dimensional manifolds whose geometry reflects latent variables tracked during sequence
processing.

Several of the clearest examples of such geometry arise in algorithmic or numerical tasks.  For instance, \citet{gurnee2025manifolds} show
that language models encode character-count variables used for line breaking on curved manifolds that are transformed across layers by
attention heads.  Similarly, recent work has shown that integer values may be represented on circular or helical manifolds that support
arithmetic operations such as modular addition \citep{kantamneni2025trigonometry}.  These results suggest that
transformers can learn geometric representations that mirror the algebraic structure of the underlying problem.

However, many of these findings concern variables with explicit numerical structure.  It remains less clear whether comparable
geometric representations arise for more abstract latent variables that arise in natural language modeling.

\subsection{Latent state tracking in transformers}

Another line of work studies how transformers track latent variables during sequence processing.  Transformers have been shown to implement
implicit Bayesian inference in certain synthetic tasks and to perform forms of meta-learning or in-context learning by representing latent
task variables in their activations \citep{akyurek2023icl, vonoswald2023transformers, xie2022meta-learning}.

In particular, recent studies have examined transformers trained on sequences generated by hidden Markov models.  In these models, the
optimal predictor must maintain a belief distribution over latent states of the generative process.  Prior work has shown that
transformers trained on such tasks can represent these belief distributions geometrically as simplices embedded in the residual
stream \citep{shai2024transformers, piotrowski2025constrained}.

These results provide compelling examples of how probabilistic beliefs over latent variables may be encoded in neural network
representations.  However, existing demonstrations largely rely on synthetic training environments where the latent variables and their
structure are known in advance.

\subsection{Discovering geometric structure in large language models}

Identifying analogous structures in large pretrained language models presents a substantially more difficult problem.  In natural language
settings, the latent variables relevant for prediction are complex and not directly observable.  As a result, geometric structures associated
with these variables cannot be located by simply aligning model representations with known ground-truth states.

Our work addresses this challenge by introducing a pipeline that searches automatically for simplex geometries in the latent feature
space of large language models.  By combining sparse feature representations with subspace clustering and geometric model fitting,
the method enables systematic discovery of candidate belief-state geometries without requiring prior knowledge of the underlying latent
variables.


% ============================================================
\section{Toy Model: Proof of Concept}
\label{sec:toy}
% ============================================================

\subsection{Generator Model}
\label{sec:toy_generator}

We construct a multipartite (MP) process by independently drawing emissions from
five component models: three Mess3 hidden Markov models and two Tom Quantum (Bloch
Walk) generalized HMMs~\citep{riechers2025quantum}.  The three Mess3 instances
have parameters $(x, a) \in \{(0.05, 0.85),\,(0.075, 0.90),\,(0.10, 0.95)\}$; the
two Tom Quantum instances have parameters $(\alpha, \beta) \in
\{(1.51, 3.07),\,(1.99, 2.51)\}$.  The MP model emission at each step is the
Cartesian product of the five component emissions, yielding a vocabulary of
$4 \times 4 \times 3 \times 3 \times 3 = 432$ tokens.

This design is \emph{intentionally challenging}.  The joint process has very high
conditional entropy: knowing the current token gives little information about the
next, because the uncertainty over each component's hidden state compounds across
five independent sources.  Random guessing achieves 0.2\% top-1 accuracy, and even
a well-trained transformer achieves only modest accuracy on this task.  We chose
this regime deliberately---not to demonstrate clean decomposition in an easy setting,
but to ask whether any signal of the underlying simplex geometries survives into the
model's representations under adverse conditions.  Prior work in controlled settings
(e.g., single-component Mess3 or HMM models) has found that transformers trained on
lower-entropy processes develop strikingly clean simplex structure~\citep{simplex_prior}.
Our goal is to probe the harder multipartite case.

We fit a small transformer to predict next tokens from this MP generator.
The architecture has $d_{\mathrm{model}} = 128$, 4 heads, 3 layers, context
length $n_{\mathrm{ctx}} = 16$, and ReLU activations (total: 709k parameters).
Training uses AdamW with base learning rate $10^{-4}$, cosine decay to 10\%,
and sequences of length 16 in batches of 2048.  The trained model achieves 7.1\%
top-1 accuracy and 24.1\% top-5 accuracy (vs.\ 0.2\% and 1.2\% at random),
confirming that the transformer learns meaningful structure while remaining far from
saturation due to the high conditional entropy of the process.

\subsection{Simplex Structure in the Residual Stream}
\label{sec:toy_pca}

\begin{figure}[t]
  \centering
  \includegraphics[width=0.75\textwidth]{figures/layer_1_pcs_3_4_TQ1.png}
  \caption{%
    \textbf{Cross-component entanglement in the toy model residual stream (layer 1,
    PCs 3--4, colored by Tom Quantum 1 sub-component output token).}
    Each point is a single generated sample.  The color of each point is the
    \emph{true output token produced by the Tom Quantum 1 sub-component} for that
    sample---one of four discrete token classes, shown in the legend.  The
    two-dimensional PCA projection (PCs 3--4) was chosen to best reveal the TQ1
    sub-component's geometry, and the four token classes are clearly separable in
    this subspace, showing that the transformer has learned to encode TQ1's
    emission in its residual stream.
    %
    Crucially, however, \emph{the same projection also reveals clear token-class
    separation for two other sub-components} (TQ2 and M33): the legend color
    clusters visible here correspond to different token classes from those
    sub-components as well, not just TQ1.  This is a direct visual signature of
    cross-component entanglement: the residual stream has not factored the five
    independent generative sources into orthogonal subspaces.  Instead, the same
    linear directions simultaneously encode the emission tokens of multiple
    sub-components.  This entanglement---not any failure of the pipeline---is the
    primary reason why unsupervised latent clustering is difficult in this setting.
    Per-component projections for all five sub-components, along with a full
    discussion of the cross-component structure, are shown in
    Appendix~\ref{app:toy_pca_clusters}.
  }
  \label{fig:toy_pca}
\end{figure}

Figure~\ref{fig:toy_pca} illustrates both the promise and the challenge of the toy
model setting.  In the PCA subspace that best reveals the Tom Quantum 1
sub-component (PCs 3--4, layer 1), the four TQ1 output token classes are cleanly
separated---demonstrating that the transformer has genuinely encoded this
sub-component's belief state in its residual stream.  At the same time, the
\emph{same} projection shows clear separation by the token classes of TQ2 and M33
as well.  The residual stream encodes multiple sub-component signals along the same
linear directions simultaneously; the five generative sources are entangled, not
factored.

This entanglement has a direct consequence for the unsupervised clustering step: any
cluster of SAE latents that tracks one sub-component's emission token will
inevitably carry correlated signal from others, making fully clean separation an
ill-posed objective.  Against this backdrop, the pipeline's ability to recover all
five sub-components with a mean $R^2$ of 0.61 (Table~\ref{tab:toy_r2},
Section~\ref{sec:toy_cluster}) represents meaningful evidence that simplex geometry
can be detected even under these adverse conditions.

\subsection{SAE Fitting and Latent Analysis}
\label{sec:toy_sae}

We fit TopK SAEs with $K \in \{3, 4, 5, 6, 7, 8, 10, 12\}$ and vanilla L1 SAEs to
each layer of the multipartite toy model ($d_{\mathrm{SAE}} = 512$).
Figure~\ref{fig:toy_epdf} shows the empirical probability density functions (EPDFs)
for all latents in the best natural cluster at layer~1, which the pipeline assigns
to the first Mess3 component (Cluster~4, $R^2 = 0.89$).  Each panel corresponds to
one of the five generative components; the assigned component is marked with a star.
Individual latents concentrate in distinct regions of each component's belief
geometry — tiling the Mess3 simplex vertices and edges, and forming localised blobs
on the Tom Quantum disk — demonstrating that the SAE partitions belief-state
information across latents in a geometry-consistent way.

\begin{figure}[t]
  \centering
  \includegraphics[width=\columnwidth]{figures/cluster_4_latent_4.png}
  \vspace{0.5em}
  \includegraphics[width=\columnwidth]{figures/cluster_4_latent_26.png}
  \caption{%
    \textbf{Two representative latents from Cluster~4 (multipartite toy model,
    layer~1, TopK $K=12$), illustrating correlated tiling across components.}
    Each panel shows the KDE-smoothed activation density for a single SAE latent
    over all five component belief geometries: equilateral triangles for the three
    Mess3 processes (belief simplex) and disks for the two Tom Quantum processes
    (equatorial Bloch sphere slice).  The assigned component is marked~$\star$
    ($R^2 = 0.89$).
    \emph{Top:} Latent~4 fires near the top vertex of each Mess3 simplex and at the
    centre of the Tom Quantum disks.
    \emph{Bottom:} Latent~26 fires near the base vertices and along the annular rim
    of the Tom Quantum disks.
    Crucially, the two latents tile complementary regions of the same geometry:
    they are correlated across components (each is coherent with the same
    ground-truth belief geometry) yet separable within each component (they
    respond to distinct belief-state regions).  This demonstrates that the pipeline
    recovers clusters whose latents partition belief-state information in a
    geometry-consistent way, even within a challenging five-component mixture.
  }
  \label{fig:toy_epdf}
\end{figure}

\subsection{Clustering and AANet Fitting}
\label{sec:toy_cluster}

We apply $k$-subspace clustering to the SAE decoder matrix to partition latents
into $N$ clusters, then fit AANet to each cluster's activation matrix with target
simplex dimensions $K \in \{2, 3, 4, 5\}$, selecting $K$ by an elbow criterion on
reconstruction loss.  For the Mess3-associated clusters the elbow consistently falls
at $K=3$, correctly identifying the underlying 2-simplex.

The central question is whether the pipeline can recover the five independent
generative components from a blind, unsupervised analysis of the SAE latent space.
To quantify this we fit a linear regression predicting each component's belief state
from the cluster latent activations and report $R^2$ on held-out data.
Table~\ref{tab:toy_r2} shows results for the best run of a hyperparameter sweep
(TopK $K=12$, $k$-subspaces, $N=6$ clusters).  Each of the five non-noise clusters
is uniquely assigned to a distinct generative component; the assignments are
one-to-one with no conflicts.

\begin{table}[t]
  \centering
  \caption{%
    \textbf{Toy model clustering recovery ($R^2$), layer 1.}
    Each cluster is assigned to the generative component whose belief state it best
    predicts (linear regression, held-out data).  All five generative components are
    recovered by a single noise-free cluster each, with no conflicts.  $R^2$ values
    are modest but uniformly positive---a nontrivial result given that the
    transformer itself achieves low predictive accuracy due to the high conditional
    entropy of the process.
  }
  \label{tab:toy_r2}
  \small
  \begin{tabular}{lcc}
    \toprule
    Cluster & Assigned component & $R^2$ \\
    \midrule
    0 & \texttt{tom\_quantum}    & 0.551 \\
    1 & \texttt{mess3\_2}        & 0.314 \\
    2 & \texttt{tom\_quantum\_1} & 0.577 \\
    3 & \texttt{mess3\_1}        & 0.719 \\
    4 & \texttt{mess3}           & 0.893 \\
    5 & (noise)                  & ---   \\
    \midrule
    \multicolumn{2}{l}{Mean (assigned clusters)} & 0.611 \\
    \bottomrule
  \end{tabular}
\end{table}

These numbers should be interpreted in light of the difficulty of the setup.
The transformer has limited ability to resolve individual component belief states,
because each token jointly encodes contributions from all five independent sources.
Matrix decomposition methods perform visibly worse here than in simpler
single-component or low-entropy settings, and the PCA projections are correspondingly
diffuse.  Against this backdrop, a mean $R^2$ of 0.61 across five independently
recovered clusters---each correctly matched to its generative component with no
supervision---represents meaningful signal.  The result supports the conclusion that
the pipeline can detect latent simplex geometry even when the underlying
representations are noisy and the predictive task is near its information-theoretic
ceiling.

% ============================================================
\section{Real Model Pipeline: Gemma-2-9B}
\label{sec:real}
% ============================================================

\subsection{SAE and Representation Space}
\label{sec:real_sae}

We use a publicly released TopK SAE ($K=50$) for Gemma-2-9B trained on the
residual stream after layer 20~\citep{saelens2024}, with $d_{\mathrm{model}} = 3584$
and $d_{\mathrm{SAE}} = 131{,}072$ latents.  Layer 20 corresponds to roughly 60\% of
the network depth, a range empirically associated with abstract contextual
representations in similar models~\citep{ferrando2024}.  We process 5000 sequences
from the Pile (context length 256) through Gemma-2-9B and store per-token SAE
activations.

\subsection{k-Subspace Clustering}
\label{sec:real_cluster}

The SAE decoder matrix $W_d \in \R^{3584 \times 131072}$ is too large to cluster in
its entirety.  We restrict to the $N$ SAE latents with the highest mean activation
frequency across the corpus, forming a candidate pool.  We cluster the corresponding
column-normalized decoder directions using k-subspace clustering with QR-pivoting
initialization.  This yields a partition of decoder directions into groups that
span approximately orthogonal subspaces of the residual stream, with the expectation
that latents in the same cluster co-activate because they jointly encode a single
underlying contextual variable.

We run the pipeline for $N \in \{512, 768\}$ (the top-$N$ most active latents by
frequency), producing $\lfloor\sqrt{N}\rfloor$ clusters per run and selecting
candidate clusters with estimated rank $k = 3$ (i.e., evidence for 2-simplex
structure from the cluster subspace elbow plot).  We additionally apply spectral
clustering as an independent check and retain only clusters identified by both
methods.

\begin{figure}[t]
  \centering
  % [FIGURE: Pipeline diagram (schematic, to be created).
  %  A horizontal flowchart showing:
  %  1. Input text → Gemma-2-9B → residual stream activations (layer 20)
  %  2. → TopK SAE → sparse latent activations
  %  3. → k-subspace clustering of decoder directions → cluster assignment
  %  4. → AANet fitting to cluster activations → barycentric coordinates
  %  5. → Validation (four strands: KL, bary>latent, causal steering, interp)
  %  This figure should be a clean schematic (e.g., TikZ) rather than a screenshot.
  %  It is the main "methods overview" figure for the paper.]
  \includegraphics[width=\textwidth]{figures/pipeline_diagram.pdf}
  \caption{%
    \textbf{Overview of the belief geometry discovery pipeline.}
    Text sequences are processed through Gemma-2-9B; residual stream activations at
    layer 20 are encoded by a TopK SAE.  Latent decoder directions are partitioned
    by k-subspace clustering, and candidate $K\!=\!3$ clusters are fit with AANet
    to recover barycentric (simplex) coordinates.  The resulting cluster
    representations are evaluated with four complementary validation measurements.
  }
  \label{fig:pipeline}
\end{figure}

\subsection{AANet Simplex Fitting}
\label{sec:real_aanet}

For each candidate cluster of size $m$ (number of latents), we form the
cluster activation matrix $\mathbf{A} \in \R^{T \times m}$ where $T$ is the
number of token positions with at least one latent active.  We project
each activation vector onto the cluster decoder subspace:
$\mathbf{x}_t = \mathbf{A}_t \mathbf{W}_c$, where $\mathbf{W}_c \in \R^{m \times
d_{\mathrm{model}}}$ is the cluster decoder submatrix.  AANet is trained on
$\{\mathbf{x}_t\}$ to minimize reconstruction loss subject to simplex
constraints, producing barycentric coordinates $\bary_t \in \Delta^{K-1}$.

The simplex dimension $K$ is selected by fitting AANet for $K \in \{2, 3, 4, 5\}$
and identifying the elbow in the reconstruction loss curve.  We focus on clusters
where the elbow is at $K = 3$ (2-simplex), since this is the smallest
non-degenerate simplex and the easiest to visualize and validate.

\subsection{Priority Cluster Selection}
\label{sec:real_selection}

\begin{table}[t]
  \centering
  \caption{%
    \textbf{The 13 priority clusters.}
    Clusters are identified by the number of top-frequency SAE latents used
    (512 or 768) and their cluster index within that run.  Confidence reflects
    manual review of interpretability outputs and validation consistency.
    $m$ denotes the number of SAE latents in the cluster.
  }
  \label{tab:priority_clusters}
  \begin{tabular}{lclll}
    \toprule
    Cluster & $m$ & Confidence & Vertex interpretation (summary) \\
    \midrule
    512\_17   & 15  & HIGH   & Text register (instructional / promotional / metadata) \\
    512\_181  & 26  & HIGH   & [Pending: see main text] \\
    768\_140  & 26  & HIGH   & Referential specificity (generic / specific / anaphoric) \\
    768\_596  & 6   & HIGH   & Grammatical person / discourse role \\
    \midrule
    512\_22   & 6   & MEDIUM & [Pending semantic label] \\
    512\_67   & 25  & MEDIUM & [Pending semantic label] \\
    512\_229  & 22  & MEDIUM & [Pending semantic label] \\
    512\_261  & 18  & MEDIUM & [Pending semantic label] \\
    512\_471  & 10  & MEDIUM & [Pending semantic label] \\
    512\_504  & 22  & MEDIUM & [Pending semantic label] \\
    768\_210  & 4   & MEDIUM & [Pending semantic label] \\
    768\_306  & 5   & MEDIUM & [Pending semantic label] \\
    768\_581  & 11  & MEDIUM & [Pending semantic label] \\
    \bottomrule
  \end{tabular}
\end{table}

After clustering and AANet fitting across both $N=512$ and $N=768$ runs, we select
13 priority clusters for validation (Table~\ref{tab:priority_clusters}).  All 13
have estimated $K=3$ from the AANet elbow criterion.  Cluster sizes range from 4 to
26 latents.  For comparison, we also select three \emph{null clusters}---random
subsets of 47--51 latents drawn without any geometric selection criterion---and
apply the identical AANet fitting and validation pipeline.  Null clusters serve as
a reference for the base rate of any measured signal under the null hypothesis of
no genuine simplex structure.

% ============================================================
\section{Validation}
\label{sec:validation}
% ============================================================

We evaluate the 13 priority clusters across four complementary validation strands.
None of the individual strands is conclusive in isolation; our argument rests on
their collective pattern.

\subsection{Strand 1: Cross-Vertex KL Divergence}
\label{sec:kl}

\paragraph{Setup.}  For each cluster, we collect near-vertex samples: token
positions where the AANet barycentric coordinate $b_k$ for one vertex exceeds a
threshold (0.70).  For each pair of near-vertex samples $(i, j)$, we run a forward
pass to obtain the model's next-token distribution $P_i, P_j \in \Delta^{V-1}$
(vocabulary size $V = 256{,}000$) and compute the symmetric KL divergence
$D_{\mathrm{KL}}(P_i \| P_j)$.  We compute three population statistics:
\emph{same-vertex} (both samples near the same vertex), \emph{cross-vertex}
(samples near different vertices), and \emph{within-simplex} (randomly sampled
active positions from the cluster, regardless of vertex proximity).

\paragraph{Hypothesis.}  If the simplex geometry is functionally meaningful---i.e.,
if different vertex positions correspond to genuinely different generative modes---then
cross-vertex pairs should have larger next-token KL divergence than same-vertex pairs:
$\overline{D}_{\mathrm{KL}}^{\mathrm{cross}} > \overline{D}_{\mathrm{KL}}^{\mathrm{same}}$.

\begin{table}[t]
  \centering
  \caption{%
    \textbf{KL divergence results: cross-vertex vs.\ same-vertex ratio.}
    We report the ratio $\overline{D}_{\mathrm{KL}}^{\mathrm{cross}} /
    \overline{D}_{\mathrm{KL}}^{\mathrm{same}}$ of mean symmetric KL divergences.
    A ratio $>1$ indicates that tokens near different simplex vertices predict more
    divergent next-token distributions than tokens near the same vertex.
    Null clusters are shown separately; their ratios overlap substantially with those
    of real clusters.
  }
  \label{tab:kl}
  \setlength{\tabcolsep}{6pt}
  \begin{tabular}{lcc|lcc}
    \toprule
    Cluster & $m$ & KL ratio & Cluster & $m$ & KL ratio \\
    \midrule
    512\_17   & 15 & 1.054 & 512\_471 & 10 & 1.017 \\
    512\_22   & 6  & 1.058 & 512\_504 & 22 & 1.164 \\
    512\_67   & 25 & 1.229 & 768\_140 & 26 & \textbf{1.987} \\
    512\_181  & 26 & 1.037 & 768\_210 & 4  & 1.017 \\
    512\_229  & 22 & 1.166 & 768\_306 & 5  & 1.102 \\
    512\_261  & 18 & 1.032 & 768\_581 & 11 & 1.144 \\
                &    &       & 768\_596 & 6  & 1.037 \\
    \midrule
    \multicolumn{6}{l}{\textit{Null clusters:}} \\
    512\_138 (null) & 47 & 1.329 & 512\_345 (null) & 51 & 1.311 \\
    768\_310 (null) & 48 & 1.108 &                 &    &       \\
    \bottomrule
  \end{tabular}
\end{table}

\paragraph{Results.} Table~\ref{tab:kl} shows that all 13 real clusters achieve
cross/same KL ratios $> 1$, ranging from 1.017 to 1.987.  However, the three null
clusters also exhibit ratios $> 1$ (1.108--1.329), sometimes exceeding those of real
clusters.  The KL divergence test thus does not cleanly discriminate real from null
clusters.  One cluster stands out markedly: \textbf{768\_140} achieves a ratio of
1.987, far higher than any other cluster (real or null), with a $z$-score of $> 4$
for the cross-vs.-same comparison.  This result is consistent with 768\_140's
strongly interpretable semantic meaning (Section~\ref{sec:highlight_140}).

\subsection{Strand 2: Latent vs.\ Barycentric Predictive Advantage}
\label{sec:bary}

\paragraph{Setup.}  The most informative test of whether the simplex geometry
captures meaningful information---beyond what any individual latent captures---is to
ask whether the full barycentric coordinate vector $\bary_t \in \R^{K}$ predicts
the next-token distribution better than the single best individual latent in the
cluster.  For each of the top-50 most frequent tokens in the near-vertex sample,
we fit a linear regression $R^2$ for the probability of that token given the
barycentric coordinate (across $T \sim 180$ samples) and separately for the
activation of each latent in the cluster.  We compare $R^2_{\mathrm{bary}}$ vs.\
$R^2_{\mathrm{best\,latent}}$ using a paired Wilcoxon signed-rank test across the 50
tokens.

\paragraph{Hypothesis.} If the simplex representation captures genuine mixture
information not present in any single latent, then $R^2_{\mathrm{bary}}$ should be
systematically larger than $R^2_{\mathrm{best\,latent}}$.

\begin{table}[t]
  \centering
  \caption{%
    \textbf{Barycentric vs.\ best-latent predictive advantage.}
    Fraction of top-50 tokens for which the barycentric coordinate achieves higher
    $R^2$ for next-token probability than the best individual latent in the cluster
    (\emph{frac. bary wins}).  Wilcoxon $p$-values test the one-sided hypothesis
    that barycentric $R^2$ is stochastically greater.  Mean $R^2$ values are shown
    for context.  Null clusters are shown at bottom; none achieves significant
    barycentric advantage.
  }
  \label{tab:bary}
  \begin{tabular}{lcccc}
    \toprule
    Cluster & Frac.\ bary wins & Wilcoxon $p$ & Mean $R^2_{\mathrm{bary}}$ & Mean $R^2_{\mathrm{best\,latent}}$ \\
    \midrule
    512\_17   & 0.00 & 1.00         & 0.055 & 0.138 \\
    512\_22   & 0.00 & ---          & ---   & --- \\
    512\_67   & 0.00 & ---          & ---   & --- \\
    \textbf{512\_181}  & \textbf{1.00} & $< 10^{-15}$ & \textbf{0.612} & 0.539 \\
    \textbf{512\_229}  & \textbf{1.00} & $< 10^{-15}$ & \textbf{0.378} & 0.244 \\
    512\_261  & 0.00 & ---          & ---   & --- \\
    512\_471  & 0.00 & ---          & ---   & --- \\
    512\_504  & 0.00 & ---          & ---   & --- \\
    768\_140  & 0.20 & $\approx 1.0$& 0.528 & \textbf{0.549} \\
    768\_210  & 0.00 & ---          & ---   & --- \\
    768\_306  & 0.16 & $\approx 1.0$& ---   & --- \\
    768\_581  & 0.00 & ---          & ---   & --- \\
    \textbf{768\_596}  & \textbf{0.98} & $< 10^{-14}$ & \textbf{0.269} & 0.217 \\
    \midrule
    512\_138 (null) & 0.00 & 1.00 & --- & --- \\
    512\_345 (null) & 0.02 & $\approx 1.0$ & --- & --- \\
    768\_310 (null) & 0.16 & $\approx 1.0$ & --- & --- \\
    \bottomrule
  \end{tabular}
\end{table}

\paragraph{Results.} Table~\ref{tab:bary} shows that 3 of the 13 real clusters
(512\_181, 512\_229, 768\_596) exhibit a significant barycentric advantage, with
barycentric $R^2$ exceeding the best individual latent's $R^2$ for nearly all or
all top-50 tokens, with Wilcoxon $p < 10^{-14}$.  Crucially, none of the three null
clusters exhibits a significant advantage ($p \approx 1.0$ in all cases).  This is
the single cleanest discriminating result in our validation battery: it directly
measures whether the \emph{mixture} information encoded in the full simplex
coordinate carries independent predictive signal, which is precisely what one would
expect if the simplex geometry reflects a latent state variable over which the model
is computing something like Bayesian mixture inference.

The improvement is substantial for 512\_181 and 512\_229 (mean $R^2$ improves by
$\sim$14\% and $\sim$55\%, respectively) but more modest for 768\_596 (improving by
$\sim$24\%).  Conversely, 768\_140---our most interpretable cluster---does not show
a significant advantage despite its strong KL divergence result; here the best single
latent ($R^2 = 0.549$) slightly exceeds the barycentric $R^2$ ($0.528$).

\subsection{Strand 3: Causal Steering}
\label{sec:steering}

\paragraph{Setup.}  We test whether the simplex geometry identified by AANet is
\emph{causally} relevant by steering the model's residual stream to target vertex
coordinates and asking whether model continuations shift to semantically match the
target vertex.  Concretely, for each source vertex $v$ (from which we start near-vertex
samples) and each target vertex $v' \neq v$, we compute an intervention direction in
residual stream space derived from the AANet vertex decoder, add it at scale
$s \in \{1, 5, 20\}$ to the residual stream at the trigger position, and generate a
continuation.  We evaluate each continuation using Qwen-2.5-72B (AWQ quantized, via
vLLM), which classifies whether the continuation matches the target vertex's semantic
description (derived from the autointerp labels; Section~\ref{sec:interp}).  The
steering score is the mean classification accuracy across source--target pairs,
types (3 intervention formulations), and scales.

A cluster \emph{qualifies} for steering analysis if at least 2 of its 3 vertices
achieve document accuracy $\geq 0.15$ (the Qwen classifier can reliably identify
natural continuation text for those vertices), ensuring the score is not trivially
driven by a degenerate vertex.  Five real clusters do not qualify; we report the
8 that do.

\begin{table}[t]
  \centering
  \caption{%
    \textbf{Causal steering scores.}
    The steering score reports the best (type, scale) combination's mean shift
    rate across all source$\to$target directions.  Clusters where AANet did not
    qualify (fewer than 2 vertices with document accuracy $\geq 0.15$) are marked
    with a dash.  Null cluster scores are shown at bottom for comparison.
    Distributions of real and null scores overlap, but the best real cluster
    (768\_596) exceeds all null clusters.
  }
  \label{tab:steering}
  \begin{tabular}{lccc}
    \toprule
    Cluster & Steering score & Best type & Best scale \\
    \midrule
    \textbf{768\_596} & \textbf{0.419} & type2 & 20 \\
    512\_22   & 0.356 & type1 & 20 \\
    768\_210  & 0.336 & type2 & 1 \\
    768\_140  & 0.289 & type1 & 20 \\
    512\_67   & 0.250 & type1 & 20 \\
    512\_261  & 0.242 & type1 & 1 \\
    768\_581  & 0.209 & type3 & 5 \\
    512\_17   & 0.118 & type3 & 5 \\
    512\_181  & ---   &  ---  & --- \\
    512\_229  & ---   &  ---  & --- \\
    512\_471  & ---   &  ---  & --- \\
    512\_504  & ---   &  ---  & --- \\
    768\_306  & ---   &  ---  & --- \\
    \midrule
    512\_138 (null) & 0.358 & type1 & 1 \\
    768\_310 (null) & 0.272 & type1 & 20 \\
    512\_345 (null) & 0.150 & type1 & 5 \\
    \bottomrule
  \end{tabular}
\end{table}

\paragraph{Results.}  Table~\ref{tab:steering} shows that 8 of 13 real clusters
qualify and achieve positive steering scores.  The strongest cluster, 768\_596,
achieves a score of 0.419.  However, the null cluster 512\_138 achieves 0.358, and
768\_310 achieves 0.272, overlapping with the lower half of the real cluster
distribution.  Causal steering thus does not cleanly discriminate real from null
clusters.

Examining the null clusters more carefully reveals a consistent pattern: all three
null clusters contain a \emph{phantom vertex} with document accuracy $< 0.15$ (i.e.,
natural continuation text from that vertex cannot be reliably classified into it).
High steering scores in these cases reflect not genuine geometric causal influence
but rather that samples labeled as originating from the phantom vertex are
misclassified in the baseline, so steering away from it yields apparent
improvement.  We note, however, that 768\_596 also has a poorly defined third vertex
(V2 document accuracy = 0.069), meaning this artifact is not unique to null clusters;
768\_596's high score should be interpreted with this caveat in mind.

Despite the overlap with null clusters, the \emph{combination} of strong steering
and barycentric advantage for 768\_596 is more compelling than either metric alone:
the barycentric advantage is not present in any null cluster, while the steering
score for 768\_596 is the highest across all clusters.

\subsection{Strand 4: Semantic Coherence and LLM Interpretability}
\label{sec:interp}

\paragraph{Setup.}  For each cluster, we collect the top-50 most activating tokens
per vertex (near-vertex samples, sorted by activation weight) and submit them for
interpretability to Claude Sonnet, repeating the labeling 20 times with different
random seeds.  We synthesize the vertex labels across iterations into a single
consolidated description and assess inter-iteration consistency.

\paragraph{Results.}  All 13 priority clusters produce semantically coherent vertex
labels that are stable across the majority of iterations.  The three best-characterized
clusters are described in detail in Section~\ref{sec:highlights}.  Semantic
coherence is necessary but not sufficient evidence: the null clusters also produce
coherent vertex labels (as expected, since any set of tokens with a common clustering
criterion will have some interpretable theme).  We treat this strand primarily as a
\emph{prerequisite check}---if we cannot produce stable, coherent labels, the other
quantitative tests lose interpretive grounding---rather than as an independent
discriminating measurement.

% ============================================================
\section{Highlight Clusters}
\label{sec:highlights}
% ============================================================

We profile three clusters that together illustrate the range of evidence available
and that, in combination, offer the most compelling case for genuine simplex
structure in Gemma-2-9B's representations.

\subsection{768\_596: Strongest Functional Evidence}
\label{sec:highlight_596}

\begin{figure}[t]
  \centering
  \begin{subfigure}[b]{0.48\textwidth}
    % [FIGURE: Latent spatial scatter for cluster 768_596.
    %  Source: outputs/validation/latent_spatial/cluster_768_596/all_latents.png
    %  This shows all 5000 sample positions in the 2D simplex (ternary plot),
    %  colored by which latent is most active (or by barycentric coordinate).
    %  The scatter should show clear triangular structure with concentration
    %  near vertices/edges, NOT a uniform cloud.]
    \includegraphics[width=\textwidth]{figures/cluster_768_596_all_latents.png}
    \caption{Latent space scatter (all samples)}
  \end{subfigure}
  \hfill
  \begin{subfigure}[b]{0.48\textwidth}
    % [FIGURE: Centroid scatter for cluster 768_596.
    %  Source: outputs/validation/latent_spatial/cluster_768_596/cluster_768_596_centroid_scatter.png
    %  This shows the simplex centroid positions of each of the 6 latents in the cluster,
    %  positioned according to their mean barycentric coordinate.
    %  The 6 centroids should cluster near the three vertices, showing that each latent
    %  has a preferred vertex position (the basis of the simplex structure).]
    \includegraphics[width=\textwidth]{figures/cluster_768_596_centroid_scatter.png}
    \caption{Per-latent centroid positions}
  \end{subfigure}
  \caption{%
    \textbf{Cluster 768\_596: spatial structure.}
    \emph{(a)} All 5000 sample positions in barycentric coordinates (ternary plot).
    Concentration near the three vertices indicates that the model visits
    near-pure vertex states, not just the simplex interior.
    \emph{(b)} Centroid positions of the 6 individual SAE latents in the cluster.
    Latents partition across the three vertices, with each vertex attracting one or two
    latents (centroids near the corresponding corner), consistent with
    vertex-specialized feature coding.
  }
  \label{fig:cluster_596_spatial}
\end{figure}

Cluster 768\_596 ($m = 6$ latents) is the single cluster with the strongest
convergence of functional evidence: it achieves the highest steering score (0.419),
and its barycentric coordinates significantly outperform the best individual latent
at next-token prediction (frac.\ bary wins $= 0.98$, $p < 10^{-14}$).
Figure~\ref{fig:cluster_596_spatial} shows the latent spatial structure: the 5000
samples occupy a triangular region with concentration near the three vertices, and
the 6 latents partition into three vertex-preferring groups.

The semantic interpretation of this cluster is moderate-confidence.  Claude Sonnet
produces stable labels across 20 iterations: V0 corresponds to third-person or
external-entity reference, V1 to first-person or self-referential framing, and V2 to
second-person or direct-address text.  However, the 6-latent cluster is small, and
the trigger token distributions at each vertex are heterogeneous, including
punctuation, function words, and content words, suggesting the true organizing
principle may be closer to \emph{syntactic discourse position} or \emph{token type
at discourse boundaries} than clean grammatical person.

\subsection{768\_140: Strongest KL Divergence and Best Interpretation}
\label{sec:highlight_140}

\begin{figure}[t]
  \centering
  \begin{subfigure}[b]{0.48\textwidth}
    % [FIGURE: Latent spatial scatter for cluster 768_140.
    %  Source: outputs/validation/latent_spatial/cluster_768_140/all_latents.png]
    \includegraphics[width=\textwidth]{figures/cluster_768_140_all_latents.png}
    \caption{Latent space scatter (all samples)}
  \end{subfigure}
  \hfill
  \begin{subfigure}[b]{0.48\textwidth}
    % [FIGURE: KL divergence distribution plot for cluster 768_140.
    %  Source: outputs/validation/kl_divergence_broad_2/768_140_kl_plot.png
    %  This shows overlaid histograms of pairwise KL divergences for:
    %  - same-vertex pairs (blue)
    %  - cross-vertex pairs (red)
    %  - within-simplex pairs (gray)
    %  The cross-vertex distribution should be clearly shifted right relative
    %  to the same-vertex distribution, with a ratio of ~2x in the means.
    %  This is the strongest KL result in the dataset.]
    \includegraphics[width=\textwidth]{figures/768_140_kl_plot.png}
    \caption{KL divergence distributions}
  \end{subfigure}
  \caption{%
    \textbf{Cluster 768\_140: spatial structure and KL divergence.}
    \emph{(a)} Sample scatter in barycentric coordinates.
    \emph{(b)} Pairwise symmetric KL divergence distributions for same-vertex
    (blue), cross-vertex (red), and within-simplex (gray) pairs.
    The cross-vertex distribution is shifted to substantially higher values,
    with a cross/same mean ratio of 1.99---by far the strongest KL signal in
    our dataset.  This indicates that tokens near different simplex vertices
    predict markedly different next-token distributions.
  }
  \label{fig:cluster_140}
\end{figure}

Cluster 768\_140 ($m = 26$ latents) achieves a KL divergence ratio of 1.987, the
largest in our dataset and far above both the other real clusters and the null clusters
(Figure~\ref{fig:cluster_140}b).  It also achieves the second-highest document
classification accuracy at vertex 1 (0.958), indicating that natural text near V1 is
highly classifiable.

The semantic interpretation of 768\_140 is the clearest of any cluster in the
dataset.  Claude Sonnet consistently assigns:
\begin{itemize}
  \item \textbf{V0}: Generic or indefinite reference---common nouns, abstract concepts,
    ``a company'', ``some idea'', ``any kind of''
  \item \textbf{V1}: Specific or definite reference---proper nouns, named entities,
    dates, ``The Apple Inc.'', ``Paris'', ``October 12''
  \item \textbf{V2}: Anaphoric or contextual reference---pronouns, demonstratives,
    ``it'', ``they'', ``this phenomenon''
\end{itemize}
This three-way partition corresponds to the classical linguistic distinction between
generic, specific, and referential noun phrase types, and maps naturally onto a
latent-state interpretation: the model may be tracking \emph{what kind of referential
context it is in}, which strongly predicts upcoming token type (common noun vs.\
proper noun vs.\ pronoun), explaining the large KL signal.

Importantly, 768\_140 does \emph{not} exhibit a significant barycentric advantage
($p \approx 1.0$; Table~\ref{tab:bary}): the best single latent outperforms
the full barycentric coordinate by a small margin.  This is consistent with a
situation where the three reference types are sufficiently discrete that a single
high-quality latent (specialized for ``proper noun'' or ``anaphoric'' reference)
already captures most of the predictive signal, leaving little room for the mixture
coordinate to add information.

\subsection{512\_17: Most Interpretable Semantic Structure}
\label{sec:highlight_17}

\begin{figure}[t]
  \centering
  \begin{subfigure}[b]{0.48\textwidth}
    % [FIGURE: Vertex latent pie charts for cluster 512_17.
    %  Source: outputs/validation/latent_spatial/cluster_512_17/cluster_512_17_vertex_latent_pies.png
    %  This shows three pie charts (one per vertex V0, V1, V2), each showing
    %  the fraction of total activation weight contributed by each of the 15
    %  SAE latents in the cluster, at tokens near that vertex.
    %  If the cluster has good simplex structure, each vertex should be dominated
    %  by different latents, with a different latent "profile" per vertex.]
    \includegraphics[width=\textwidth]{figures/cluster_512_17_vertex_latent_pies.png}
    \caption{Per-vertex latent activation profiles}
  \end{subfigure}
  \hfill
  \begin{subfigure}[b]{0.48\textwidth}
    % [FIGURE: Centroid scatter for cluster 512_17.
    %  Source: outputs/validation/latent_spatial/cluster_512_17/cluster_512_17_centroid_scatter.png]
    \includegraphics[width=\textwidth]{figures/cluster_512_17_centroid_scatter.png}
    \caption{Per-latent centroid positions}
  \end{subfigure}
  \caption{%
    \textbf{Cluster 512\_17: latent decomposition.}
    \emph{(a)} Fraction of total cluster activation contributed by each latent at
    each of the three vertices.  Different latents dominate at each vertex, consistent
    with vertex-specialized feature encoding.
    \emph{(b)} Centroid positions of the 15 SAE latents.  Clear separation across
    the three simplex vertices confirms that the cluster's latents have distinct
    vertex preferences.
  }
  \label{fig:cluster_17}
\end{figure}

Cluster 512\_17 ($m = 15$ latents, HIGH confidence) achieves the clearest and most
consistent semantic interpretation of any cluster.  Claude Sonnet assigns, with
high stability across 20 iterations:
\begin{itemize}
  \item \textbf{V0} [consistent]: Instructional/procedural text---how-to guides,
    technical step descriptions, user-action sequences
  \item \textbf{V1} [consistent]: Descriptive/promotional text---service descriptions,
    product features, business offerings
  \item \textbf{V2} [mixed]: Metadata/navigational/organizational context---company
    information, navigation elements, structural text
\end{itemize}
This partition corresponds to \emph{document register} or text genre: the model may
be tracking what type of discourse mode it is in, which governs the expected vocabulary
and syntax of subsequent text.  The belief-state interpretation is natural: the
``generative process'' is the human author, and the hidden state is the register they
are writing in.

Quantitatively, 512\_17 is one of the weaker clusters: its steering score (0.118) is
the lowest among qualified clusters, and it shows no significant barycentric advantage.
This may reflect that register is more distributed across many SAE features and that
the 15-latent cluster captures only a portion of the relevant representational
subspace, making causal interventions from it relatively weak.

% ============================================================
\section{Discussion}
\label{sec:discussion}
% ============================================================

\subsection{Convergent Evidence Without a Single Decisive Test}
\label{sec:convergent}

No single validation measure cleanly separates all real clusters from null clusters.
The KL divergence ratio is elevated for real clusters overall, but null clusters
achieve comparable ratios.  Causal steering is positive for most qualified real
clusters, but one null cluster (512\_138, 0.358) is among the strongest in the
dataset.  The one measure that does cleanly discriminate is the barycentric predictive
advantage: 3 real clusters achieve significant advantages ($p < 10^{-14}$) while
zero null clusters do.

The fragmentation of evidence across clusters---768\_596 has the best
steering and barycentric advantage but murky semantics; 768\_140 has the best KL and
interpretability but no barycentric advantage; 512\_17 has the clearest semantics
but weak quantitative validation---is itself informative.  It suggests that the
discovered structures are not all the same kind of object: some clusters may track
discrete, interpretable distinctions (768\_140's reference types) that are already
captured by a single strong latent, while others may encode graded mixture states
(768\_596) where the coordinate information provides genuine marginal signal.

\subsection{Comparison to Toy Model Expectations}
\label{sec:toy_comparison}

The toy model establishes that, in a fully controlled setting where the generative
process is known, the pipeline correctly recovers the ground-truth simplex geometry
from SAE latent space with imperfect but directionally correct reconstruction.  The
real-model results are qualitatively consistent with the toy model: we find
simplex-like structures that are partially functional, semantically coherent, and
causally relevant but not as cleanly recoverable as in the controlled case.  The
toy model results contextualize this: even with perfect ground truth, AANet
reconstruction is noisy, so the real-model results are consistent with genuine but
imperfect encoding of belief geometries under considerably more complex generative
dynamics.

\subsection{Limitations}
\label{sec:limitations}

\paragraph{Effect sizes are modest.}
The best steering score (0.419) is well below 1.0.  The barycentric $R^2$ advantage,
while significant, corresponds to modest absolute improvements.  This is consistent
with the simplex structures being present but fragile or partially encoded.

\paragraph{Null cluster size mismatch.}
Our null clusters contain 47--51 latents compared to 4--26 for real clusters.
Larger clusters provide more activation signal, potentially inflating baseline
measures.  A more controlled comparison would match null clusters to real cluster sizes.

\paragraph{Phantom vertices.}
Several clusters---real and null---have one poorly-defined vertex with document
classification accuracy below 0.15.  Steering and classification results involving
phantom vertices are unreliable.  We note and exclude these where possible but cannot
fully eliminate the confound.

\paragraph{No causal proof of belief-state tracking.}
We cannot directly observe what internal computation Gemma-2-9B performs, and our
evidence is correlational (KL, barycentric $R^2$) or weakly causal (steering, which
is a perturbation experiment but not a controlled identification of causal
mechanisms).  The findings are consistent with belief-state tracking but do not rule
out alternative explanations.

\paragraph{Single model, single layer.}
All real-model results are from Gemma-2-9B, layer 20.  Whether the findings
generalize to other models, architectures, or layers is unknown.

\paragraph{Degeneracy at large steering scales.}
At scale $s = 20$, many clusters produce high-degeneracy continuations (model
outputs non-text), indicating the model's representations are disrupted before
exhibiting clean state transitions.  This bounds the strength of causal evidence
we can obtain.

% ============================================================
\section{Conclusion}
\label{sec:conclusion}
% ============================================================

We have introduced and validated a pipeline for finding simplex-structured
representational geometries in SAE latent spaces of language models, motivated by the
hypothesis that language models track latent generative states as internal belief
states.  In a controlled toy model setting, the pipeline correctly recovers known
simplex structures from the SAE latent space.  Applied to Gemma-2-9B, we identify
13 priority clusters exhibiting $K=3$ simplex structure, of which 3 show a
significant predictive advantage of barycentric over individual latent coordinates---a
signal not present in null clusters.  Causal steering confirms that AANet-derived
simplex coordinates are causally connected to model behavior in 8 of 13 clusters,
and all clusters produce semantically coherent and stable vertex labels.

The strongest individual cluster, 768\_140, represents the linguistic distinction
between generic, specific, and anaphoric reference---a semantically natural candidate
for a latent generative state tracked across context.  The functionally strongest
cluster, 768\_596, exhibits both barycentric predictive advantage and the highest
causal steering score in the dataset.  Together, these findings suggest that
Gemma-2-9B's SAE latent space contains subspaces that, at a minimum, encode
mixture-like representations of functionally distinct discourse modes---representations
that behave in ways expected of belief-state encodings, even if their precise
computational role remains to be established.

Future work should address the limitations identified above, including better-matched
null clusters, temporal analysis of barycentric coordinates across context windows
(to test whether they update in ways consistent with Bayesian belief updating), and
extension to additional models and layers.

% ============================================================
% References
% ============================================================

\bibliographystyle{plainnat}
\bibliography{references}

% ============================================================
% Appendix
% ============================================================
\appendix

\section{Toy Model Details}
\label{app:toy}

\paragraph{Mess3 process.}  A Mess3 HMM~\citep{marzen2017predictive} has three
hidden states $\{0, 1, 2\}$ with transition probabilities parameterized by
$(x, a)$: the self-transition probability is $a$, and the remaining probability
$1-a$ is split uniformly among the other two states.  The emission from state $s$
is token $s$ with probability $x$ and a random token with probability $1-x$.

\paragraph{Tom Quantum (Bloch Walk) process.}  The Bloch Walk GHMM~\citep{riechers2025quantum}
is parameterized by $(\alpha, \beta)$ controlling the step distribution on the Bloch
sphere.  The internal state is a unit vector on $S^2$; emissions are four tokens
corresponding to projective measurements.  The resulting temporal correlations
exhibit non-Markovian structure.

\paragraph{Multipartite model vocabulary.}  The MP model generates tuples from
$\{0,\ldots,3\}^2 \times \{0,1,2\}^3$ (two Bloch Walk components $\times$ three
Mess3 components), encoded as a single integer token from $\{0,\ldots,431\}$.

\section{Toy Model: Per-Component PCA Projections}
\label{app:toy_pca_clusters}

Figure~\ref{fig:app_toy_cluster_pca} shows the best PCA projections we found for
each of the five generative components at layer 1, colored by ground-truth belief
state.  Two features of these plots are worth noting.

First, different components are best revealed by different principal components: the
two Tom Quantum components appear in PCs 1--4, while the three Mess3 components
require PCs 5--11.  This confirms that each component has carved out a partially
distinct subspace in the residual stream.

Second, and more strikingly, the legend labels in each panel are the true tokens
emitted by each component---that is, the discrete class label from that source
for each generated sample.  Each panel is therefore colored by the emitted token
class of the component it is designed to show, but it also depicts the emitted
tokens of the other components as visual overlays.  What we find is that each plot
reveals clear token-class separation for \emph{other} components beyond the one it
was designed to show.  The TQ2 projection (PCs 1--2) shows clean separation by TQ1
token class as well.  The TQ1 projection (PCs 3--4) reveals separation for both TQ2
and M33.  The M33 projection shows clear separation for TQ1 and M32.  Most
strikingly, the M31 and M32 projections (PCs 8--11) both show near-perfect
token-class separation for \emph{each other} simultaneously.

This is not merely a matter of sharing PC indices: it indicates that the residual
stream encodes the emission token of multiple components along the same linear
directions.  The transformer has not factored the five components into orthogonal
subspaces; their signals are genuinely entangled.  This is a direct consequence of
the joint multipartite token structure, and it is what makes the unsupervised
clustering problem hard: any cluster whose activations track one component's token
will inevitably carry correlated signal from others.  Despite this, the recovered
clusters achieve a mean $R^2$ of 0.61 (Table~\ref{tab:toy_r2}).

\begin{figure}[p]
  \centering
  \begin{subfigure}[t]{0.48\textwidth}
    \centering
    \includegraphics[width=\textwidth]{figures/layer_1_pcs_1_2_TQ2.png}
    \caption{%
      \textbf{Tom Quantum 2, PCs 1--2.}
      Each point is colored by the true output token of the \emph{TQ2}
      sub-component (4 classes).  TQ2's token classes are fully separated in this
      2D projection, tracing an arc consistent with Bloch-sphere geometry.
      Notably, TQ1's token classes are also clearly separated in this same
      subspace, indicating that PCs 1--2 encode both Tom Quantum sub-components
      simultaneously.
    }
    \label{fig:app_pca_tq2}
  \end{subfigure}
  \hfill
  \begin{subfigure}[t]{0.48\textwidth}
    \centering
    \includegraphics[width=\textwidth]{figures/layer_1_pcs_3_4_TQ1.png}
    \caption{%
      \textbf{Tom Quantum 1, PCs 3--4.}
      Each point is colored by the true output token of the \emph{TQ1}
      sub-component (4 classes).  TQ1's token classes separate cleanly into
      arc/circular structure in this 2D projection.  However, the same subspace
      also reveals clear token-class separation for TQ2 and M33, making this
      panel the most striking illustration of cross-component entanglement: three
      independent sub-components are simultaneously legible in a single 2D
      projection.  (Also shown as Figure~\ref{fig:toy_pca} in the main text.)
    }
    \label{fig:app_pca_tq1}
  \end{subfigure}

  \vspace{1em}

  \begin{subfigure}[t]{0.32\textwidth}
    \centering
    \includegraphics[width=\textwidth]{figures/layer_1_pcs_5_6_7_M33.png}
    \caption{%
      \textbf{Mess3 \#3, PCs 5--6--7.}
      Each point is colored by the true output token of the \emph{M33}
      sub-component (3 classes).  The 3D projection reveals triangular
      (2-simplex) structure for M33.  Token-class separation for TQ1 and
      M32 is also visible in this subspace.
    }
    \label{fig:app_pca_m33}
  \end{subfigure}
  \hfill
  \begin{subfigure}[t]{0.32\textwidth}
    \centering
    \includegraphics[width=\textwidth]{figures/layer_1_pcs_8_9_11_M31.png}
    \caption{%
      \textbf{Mess3 \#1, PCs 8--9--11.}
      Each point is colored by the true output token of the \emph{M31}
      sub-component (3 classes).  M31's 3 token classes are clearly
      separated in triangular geometry.  M32's token classes are also
      near-perfectly separated in this same subspace (see
      Figure~\ref{fig:app_pca_m32}), as PCs 8--11 encode both Mess3
      sub-components simultaneously.
    }
    \label{fig:app_pca_m31}
  \end{subfigure}
  \hfill
  \begin{subfigure}[t]{0.32\textwidth}
    \centering
    \includegraphics[width=\textwidth]{figures/layer_1_pcs_9_10_11_M32.png}
    \caption{%
      \textbf{Mess3 \#2, PCs 9--10--11.}
      Each point is colored by the true output token of the \emph{M32}
      sub-component (3 classes).  M32's triangular geometry is clear, and
      M31's token classes are simultaneously near-perfectly separated in
      this subspace---the clearest example in the dataset of two distinct
      sub-components sharing an almost identical representational subspace.
    }
    \label{fig:app_pca_m32}
  \end{subfigure}

  \caption{%
    \textbf{Per-sub-component PCA projections, toy model layer 1.}
    Each panel shows the 2D or 3D PCA subspace that best reveals one
    sub-component's geometry.  \emph{In every panel, each point is colored by the
    true discrete output token produced by that panel's target sub-component for
    that sample}---not by a continuous belief state or any derived quantity.  The
    legend shows the token class labels.  The target sub-component's token classes
    are cleanly separable in each panel's subspace.  However, every panel also
    exhibits clear token-class separation for at least one other sub-component,
    demonstrating that the residual stream encodes multiple independent generative
    sources along the same linear directions.  This cross-component entanglement
    is the primary source of difficulty for unsupervised latent clustering in this
    setting.
  }
  \label{fig:app_toy_cluster_pca}
\end{figure}

\section{SAE Details}
\label{app:sae}

We use the Gemma-2-9B SAE from the SAELens library (release
\texttt{gemma-scope-9b-pt-res}, layer 20, expansion factor $\approx 36.6$, $K = 50$
active features per token).  The SAE was trained by DeepMind on a proprietary mixture
of text data using TopK activation with $L_2$ reconstruction loss; full details are
described in the associated technical report.

For the toy model SAEs, we use a custom TopK SAE implementation with $d_{\mathrm{SAE}}
= 256$, training for 20{,}000 steps on sequences generated online from the MP model.
We tried $K \in \{3, 4, 5, 6, 7\}$; the best-recovering SAE for the Mess3 component
uses $K = 3$.

\section{AANet Architecture and Training}
\label{app:aanet}

We use the AANet implementation from \citet{korem2019finding}.  The encoder consists
of two hidden layers (widths 256, 128 by default), a bottleneck of dimension $K-1$
(to enforce the simplex constraint via softmax normalization), and an analogous
decoder.  The loss combines $L_2$ reconstruction, a simplex penalty
$\lambda_1 (1 - \|\mathbf{E}(a)\|_1)$, and a non-negativity penalty
$\lambda_2 \sum_i |x_i| \mathbf{1}(x_i < 0)$.  We train with AdamW,
learning rate $10^{-3}$, for 10{,}000 steps, selecting the checkpoint with the
lowest validation reconstruction loss.  We run 5 independent restarts and average
the barycentric coordinates.

\section{Causal Steering Intervention Types}
\label{app:steering_types}

We use three types of steering interventions, applied at the residual stream at the
trigger token position:

\begin{itemize}
  \item \textbf{Type 1}: Ablate and replace.  Remove the current cluster activation
    contribution (subtract the projection onto the cluster subspace) and add back the
    cluster activation corresponding to the target vertex decoder output.
  \item \textbf{Type 2}: Add direction.  Add the difference of AANet decoder outputs
    at target vs.\ source vertex barycentric coordinates, scaled by $s$.
  \item \textbf{Type 3}: Add decoder direction.  Add the raw AANet decoder output
    for the target vertex barycentric coordinate, scaled by $s$, without subtracting
    the source vertex contribution.
\end{itemize}

\section{All Priority Cluster Results}
\label{app:all_results}

\begin{table}[h]
  \centering
  \caption{%
    \textbf{Full results across all four validation strands for the 13 priority clusters
    and 3 null clusters.}
    KL ratio: cross/same symmetric KL divergence ratio.
    Bary wins: fraction of top-50 tokens where $R^2_{\mathrm{bary}} > R^2_{\mathrm{best\,latent}}$
    (``sig.'' if Wilcoxon $p < 0.001$).
    Steering: best-combination steering score (``---'' if cluster did not qualify).
    Interp.: consistency of frontier-LLM vertex labels across 20 iterations
    (HIGH/MED/LOW).
  }
  \label{tab:full_results}
  \begin{tabular}{lccccc}
    \toprule
    Cluster & $m$ & KL ratio & Bary wins & Steering & Interp. \\
    \midrule
    512\_17   & 15 & 1.054 & 0.00       & 0.118 & HIGH \\
    512\_22   & 6  & 1.058 & 0.00       & 0.356 & MED \\
    512\_67   & 25 & 1.229 & 0.00       & 0.250 & MED \\
    512\_181  & 26 & 1.037 & 1.00 (sig.)& ---   & HIGH \\
    512\_229  & 22 & 1.166 & 1.00 (sig.)& ---   & MED \\
    512\_261  & 18 & 1.032 & 0.00       & 0.242 & MED \\
    512\_471  & 10 & 1.017 & 0.00       & ---   & MED \\
    512\_504  & 22 & 1.164 & 0.00       & ---   & MED \\
    768\_140  & 26 & \textbf{1.987} & 0.20 & 0.289 & MED \\
    768\_210  & 4  & 1.017 & 0.00       & 0.336 & MED \\
    768\_306  & 5  & 1.102 & 0.16       & ---   & MED \\
    768\_581  & 11 & 1.144 & 0.00       & 0.209 & MED \\
    768\_596  & 6  & 1.037 & 0.98 (sig.)& \textbf{0.419} & MED \\
    \midrule
    512\_138 (null) & 47 & 1.329 & 0.00 & 0.358 & --- \\
    512\_345 (null) & 51 & 1.311 & 0.02 & 0.150 & --- \\
    768\_310 (null) & 48 & 1.108 & 0.16 & 0.272 & --- \\
    \bottomrule
  \end{tabular}
\end{table}

\end{document}
